\section{Direct evaluation does not justify the abstention rule}
\label{sec:calibration}

\paragraph{The claim being re-evaluated.}
The project gate memo, Section 15, states: ``The pre-committed cutpoint is
validated on real cells.'' Its evidence is recovery between 0.85 and 1.08
above the 50-cell reporting boundary and 1.18 in the 20--49-cell band, based
on six draws per setting at 1{,}200 cells per arm. That comparison is an
observed/target point-estimate ratio. In contrast, the design specification,
Section 4, requires at least 90\% interval inclusion and 80\% exclusion of zero
at a halving of expression. These passages document the mismatch within this
project; they do not establish a general practice. Appendix~\ref{app:provenance}
provides the source locations and preserves their status as an internal claim.

\paragraph{Target, interval and operating requirements.}
The design document fixed a detectable halving of per-cell output ($s=0.5$),
interval inclusion of the stored reference in at least $90\%$ of defined
outputs, and discrimination of at least $80\%$, defined as the probability
that the interval excludes zero at that effect. We call the first rate benchmark inclusion (``coverage'' in the saved tables).
The candidate $ok$ is the first evaluated count or bin meeting both criteria;
$wide$ is the first meeting coverage alone. The committed values are 50 and 20.
Candidates do not require every larger count to qualify and are not certified
thresholds.

The calibration uses paired subsets of the published colorectal atlas
\citep{lee2020}: ten patients from GSE132465 (SMC) and six from
GSE144735 (KUL3). We analyse GUCA2A in the differentiated population under fixed
lineage labelling. Each draw samples, with replacement, 2{,}000 cells per arm
from an eligible two-patient pool; the reference mature fraction is $f_N=0.40$.
The implemented pool combines cells from both tissues before drawing either
synthetic arm. Because the design document did not specify tissue eligibility,
we also evaluate a reference-tissue-only pool. Neither reconstructs the real
tissue contrast. A nominal 95\% percentile interval resamples mature cells in
both arms, with fractions held fixed and 200 inner draws unless stated otherwise.
Undefined intervals are excluded from rate denominators and reported as
abstentions (Appendix~\ref{app:gridcaveats}).

The stored ``parametric'' intrinsic reference is
\begin{equation}
i_{\mathrm{req}}(D_N;s)=f_N(s-1)\bar Y_N,
\label{eq:requested}
\end{equation}
where $\bar Y_N$ is the mature-cell mean of the \emph{drawn} reference arm.
Thus the shift is fixed but this benchmark varies across draws; it is not the
fixed source-pool parameter $f_N(s-1)\mu_{\mathrm{pool}}$.
Benchmark inclusion below uses Equation~\ref{eq:requested}. This is a legitimate
draw-specific performance question, but it differs from fixed-population
coverage. We preserve that specification to test the original claim; a
population-coverage study requires a separately fixed target and sampling law.

\paragraph{Controlled comparison of counts, bins and source pools.}
Historical grids returned reference-only binned candidates of 100, 90 and 70,
but changed random streams, bin boundaries and replicate budgets together
(Appendix~\ref{app:historical}). To isolate these choices, we simulated the
union of counts once per cohort, pool and each of eight seeds, with 200 outer
replicates and 200 inner bootstrap draws. Every grid summary reuses the same
outputs at shared counts. The committed grid reaches
$\{0,20,40,100,200,400,800\}$; filled and dense grids add count support.
Fixed bins share boundaries across grids; adaptive bins recalculate them.
Binned candidates are median counts within the selected bin, not estimated
crossing locations.

\begin{figure}[htbp]
\centering
\includegraphics[width=\linewidth]{figures/controlled_calibration.pdf}
\caption{Direct evaluation on the controlled union of counts, without binning.
Rows show cohorts; columns show draw pools. Curves are median rates over eight
seeds and bands span their minima and maxima, conditional on the observed
cohorts. Dashed lines mark the 90\% benchmark-inclusion and 80\% discrimination targets;
the vertical line marks the committed count of 50. Rates condition on defined
intervals and coverage uses the draw-dependent reference in
Equation~\ref{eq:requested}. Zero-cell settings have no defined interval and
are omitted. No pooled-cell candidate qualifies. Per-count reference-only
candidates vary by seed and cohort (Table~\ref{tab:controlledgrid}).}
\label{fig:calibration}
\end{figure}

With common fixed bins, SMC returns 100 on the committed grid and 70 on both
expanded grids in all eight seeds. Adaptive bins return 65 on seven filled-grid
seeds and 90 on one, while the dense grid returns 70 throughout. KUL3 instead
returns 400 on six seeds and 800 on two, with all grid and binning schemes
agreeing within each seed. Pooled draws yield no joint candidate in either
cohort. The largest binomial Monte Carlo standard error is 0.035, conditional
on the fixed cohorts; it does not measure patient-sampling uncertainty.

\begin{table}[htbp]
\centering\small
\caption{Joint coverage-and-discrimination candidates from the controlled
union-grid run. Entries give candidates with seed frequencies in parentheses
where they vary; otherwise all eight seeds agree. Every grid summary reuses the
same draws at shared counts. ``None'' means no tested count or bin qualifies.}
\label{tab:controlledgrid}
\begin{tabular}{@{}llccc@{}}
\toprule
cohort / pool & grid & fixed bins & adaptive bins & per count \\
\midrule
SMC / reference & committed & $100$ & $100$ & $100$ \\
 & filled & $70$ & $65\ (7),\ 90\ (1)$ & $50$--$80$ \\
 & dense & $70$ & $70$ & $45$--$80$ \\
KUL3 / reference & all three & $400\ (6),\ 800\ (2)$ & same & same \\
Both / pooled & all three & None & None & None \\
\bottomrule
\end{tabular}
\end{table}

At the committed count of 50, SMC reference-only discrimination ranges
from 75.5\% to 81.0\% across seeds, with both criteria met on two of eight
seeds. KUL3 reference-only discrimination is only 45.0--55.0\% at that count.
These direct observations make the lack of support for a cohort-independent
50-cell rule more concrete than the selected binned candidates alone.

The per-count summaries in Table~\ref{tab:controlledgrid} supply a useful
check on the bin medians. In SMC, the filled grid's first qualifying actual
count is 50 on two seeds, 60 on five and 80 on one; the dense grid adds a
candidate of 45 on one seed. Thus a stable binned candidate of 70 is not
evidence that 70 is a uniquely determined biological threshold. A bin pools
outcomes over several count settings, whereas a per-count rate addresses only
that setting. Fixed bins hold the boundaries constant, but adding settings can
still change both a bin's median count and its aggregate performance.

\paragraph{Low-count rejection and nonmonotone qualification.}\label{sec:lowcounts}
The high exclusion rate at very small counts in Figure~\ref{fig:calibration}
is not evidence of reliable low-count power. We performed an exploratory
diagnostic at counts 5, 50, 100 and 800: 200 replicates per cohort, pool and
shift ($s=0.5$ and $1$), using the first controlled-run seed and the same
200-resample interval. All 16 alternative-setting summaries reproduce their
saved counterparts. Under the null ($s=1$), exclusion of zero is a false
positive, and the target is exactly zero under both the draw-dependent and
fixed-pool definitions. Table~\ref{tab:lowcounts} reports both shifts with the
same method and valid-pair denominator; all 6{,}400 diagnostic outputs are finite.

\begin{table}[htbp]
\centering\small
\caption{Exclusion of zero (\%) with the same percentile interval under a halving
(effect, $s=0.5$) and the null ($s=1$). Each entry uses 200 finite draws from
one seed stream; nominal null rejection is 5\%. Diagnostic runs are exploratory.
Five-cell intervals are evaluated before the reporting gate, which abstains
below 20 cells.}
\label{tab:lowcounts}
\begin{tabular}{@{}lrrrrrrrr@{}}
\toprule
& \multicolumn{2}{c}{5 cells} & \multicolumn{2}{c}{50 cells} & \multicolumn{2}{c}{100 cells} & \multicolumn{2}{c}{800 cells} \\
cohort / pool & effect & null & effect & null & effect & null & effect & null \\
\midrule
SMC / pooled & 70.5 & 63.0 & 52.5 & 28.5 & 56.0 & 18.0 & 69.5 & 9.5 \\
SMC / reference & 58.0 & 37.5 & 78.0 & 12.0 & 88.5 & 7.0 & 98.5 & 7.0 \\
KUL3 / pooled & 75.5 & 67.0 & 54.0 & 30.0 & 33.5 & 19.0 & 53.5 & 4.5 \\
KUL3 / reference & 70.0 & 61.5 & 55.0 & 15.5 & 56.0 & 16.0 & 87.0 & 6.5 \\
\bottomrule
\end{tabular}
\end{table}


At five cells, null rejection is 37.5--67.0\%, far above the nominal 5\%.
For the pooled SMC and KUL3 alternatives, respectively 96 and 85 of 200 tumour
samples have all-zero target expression; all of these exclude zero. Resampling
such a tumour sample cannot recover its unobserved positive tail. The reference
arm still varies, so none of the 6{,}400 full intervals has exactly zero width.
The pooled five-cell median widths are 1.45 and 1.72, compared with 2.35 and
2.01 at 50 cells. These checks support sparse-sample undercoverage, rather
than a claim that the whole interval collapses to a point. Even at 50 cells,
null rejection is 12.0--30.0\%. The original 90\% benchmark-inclusion criterion
at $s=0.5$ therefore cannot substitute for checking the null.

Nonmonotonicity also changes candidate interpretation. In the full saved
union grid, three of eight SMC/reference seeds fail a criterion at a larger
count after first qualifying. Requiring every subsequently evaluated count
to qualify changes first candidates 45, 50 and 50 to 60, 70 and 70. No such
reversal occurs in the eight KUL3/reference seeds; neither pooled condition
has a qualifying candidate. This is a descriptive re-reading of the same
simulations, not an independently validated replacement selection rule.
Together, the null diagnostic and reversals make the candidates useful as
sensitivity summaries, not operational cutpoints. The diagnostic uses one
seed stream and fixed observed cohorts; its binomial Monte Carlo standard
errors are at most 3.54 percentage points, excluding patient uncertainty.

\paragraph{Simulation precision is not patient-level uncertainty.}
Enumerating all 45 and 15 two-patient holdout pairs separates patient omission
from repeated simulation. Keying streams to the sorted pair identifiers keeps
every retained pair fixed. Raising the inner bootstrap budget from 200 to 1{,}000
and 5{,}000 leaves the first cohort's candidate at 70, but moves the second's
400 to 300 and back to 400. Omitting one source patient changes the joint
coverage-and-discrimination candidate for none of ten patients in the first
cohort and three of six in the second (KUL01, KUL30 and KUL31). For the
coverage-only candidate, the corresponding counts are five of ten and zero of
six. These are cohort-conditional influence checks, not intervals for a larger
patient population. We therefore report candidates, not an endorsed calibrated
cutpoint. Appendix~\ref{app:guards} records further biological comparability limits.

\begin{table}[htbp]
\centering\small
\caption{Reference-only sensitivity with every two-patient holdout pair
represented once per count and effect. The three inner budgets reuse identical
outer samples. Omission counts compare with the baseline at 200 inner draws;
all pairs involving the omitted patient are removed and retained pairs are
unchanged. These are sensitivity summaries within the observed cohorts.}
\label{tab:sensitivity}
\begin{tabular}{@{}lcrrrcc@{}}
\toprule
 & pairs & \multicolumn{3}{c}{joint candidate, inner draws} &
 \multicolumn{2}{c}{omissions changing candidate} \\
\cmidrule(lr){3-5}\cmidrule(lr){6-7}
cohort & & $200$ & $1{,}000$ & $5{,}000$ & joint & coverage only \\
\midrule
SMC & $45$ & $70$ & $70$ & $70$ & $0/10$ & $5/10$ \\
KUL3 & $15$ & $400$ & $300$ & $400$ & $3/6$ & $0/6$ \\
\bottomrule
\end{tabular}
\end{table}

The sensitivity windows cover counts $40,50,60,70,80,90,100,110$ in SMC and
$200,300,400,600,800$ in KUL3. These per-count summaries use a fixed pair
schedule, so they are a separate experiment from the eight-seed union grid.
Omitting KUL01 or KUL30 moves the joint candidate from 400 to 300; omitting
KUL31 moves it to the lowest tested count, 200, with its lower boundary
unobserved. The coverage-only baselines likewise occur at the lowest tested
counts, 40 and 200. Their stability cannot locate a lower crossing outside
those windows. Enumerating pairs improves transparency about source-patient
influence, but overlapping pairs do not create additional independent patients.


\paragraph{Biological controls fail to establish mechanism.}\label{sec:biocontrols}
The pre-committed panel included compositional, intrinsic and retained-gene
controls. In the primary cohort the retained-gene control changed by
$-1.000$ to $-0.951$ relative to baseline, almost matching the compositional
controls' $-1.000$ to $-0.952$. The intrinsic control had only 0.04 CP10K,
compared with 19.4 and 24.9 for the abundant compositional controls (counts per
10{,}000), making it uninformative for the intended test. A second cohort
reproduced the pattern. Thus the application has two distinct unresolved
problems: its interval rule lacks the recorded calibration support, and its
control panel does not distinguish the proposed biological mechanisms.
Improving one would not establish the other.
